\chapter{The Jacobian as an abelian variety}
\label{chap:Jacobian}
Let \(C\) be a smooth proper curve over \(k\) in the sense of Chapter~\ref{chap:Genus}, with genus \(g = g(C)\).
The Jacobian is built \emph{uniformly} as the Picard identity component \(J = \Pic^0_{C/k}\), an abelian variety of dimension \(\genus C\), produced by \textbf{Route~A (the Picard scheme via FGA representability)}. There is no genus stratification: the same construction applies for every genus, and the genus-\(0\) case is \emph{not} special --- when \(\genus C = 0\) one has \(H^1(C, \mathcal O_C) = 0\), so \(\Pic^0_{C/k}\) is automatically the trivial abelian variety \(\Spec k\). The construction is the project's committed critical path: the object must be a genuine \(\genus C\)-dimensional abelian variety even when \(C(k) = \emptyset\) (see \cref{thm:nonempty_jacobianWitness} Route~A). The FGA representability for \(\Pic_{C/k}\) and its Hilbert/Quot prerequisites are the single largest Mathlib build-out the project gates on. The Albanese universal property of \(\Pic^0_{C/k}\) consumes the general Milne \S I.1--I.3 rigidity-lemma corollaries of \cref{chap:AbelianVarietyRigidity}.
\section{The Albanese construction}
For a smooth proper curve \(C\) over \(k\), the Albanese variety is the universal abelian variety receiving a morphism from \(C\). Over an algebraically closed field this is classically identified with the connected component of the identity of the Picard scheme, but we construct it directly from symmetric powers.
\subsection{Symmetric powers and the Abel--Jacobi map}
Let \(C^{(g)} = \Sym^g(C)\) be the \(g\)\textsuperscript{th} symmetric power of \(C\) (the quotient of \(C^g\) by the action of the symmetric group \(S_g\)). It is a smooth proper scheme over \(k\) of dimension \(g\). Fix a \(k\)-rational point \(P_0 \in C(k)\). The map
\[
  \sigma \colon C^{(g)} \longrightarrow \Pic^g_{C/k}, \qquad
  \sum_{i=1}^g [P_i] \longmapsto \mathcal O_C\bigl(\textstyle\sum_{i=1}^g P_i - g P_0\bigr)
\]
is surjective and birational onto its image; its fibres are projective spaces (Brill--Noether--Riemann--Roch). Choosing \(P_0\) gives a translation isomorphism \(\Pic^g_{C/k} \xrightarrow{\;\sim\;} \Pic^0_{C/k}\), and the Stein factorisation of \(\sigma\) yields the Jacobian as the connected component of the identity.
In the Lean formalisation we do not have general symmetric powers of schemes (or quotients by finite group actions) in Mathlib. We therefore take a more direct approach: we define the Jacobian as the Albanese object of \(C\) inside the category of smooth proper group schemes over \(k\), using the universal property and the known dimension formula.
\subsection{The Albanese universal property}
\begin{definition}\leanok
[Albanese universal property of a pointed curve]
  \label{def:IsAlbanese}
  \lean{AlgebraicGeometry.IsAlbanese}
  % SOURCE: [Kleiman, The Picard scheme], Remark rmk:Alb (the Albanese map), FGA Explained ch.~5; arXiv:math/0504020, p.~50  (read from references/kleiman-picard-src/kleiman-picard.tex)
  % SOURCE QUOTE: "Reversing the preceding argument, we see that every
  % such \(b\) arises from a unique map \(a\:B^*\to \IPic_{X/k}\) such that
  % \(a(0)=0\).  Since \(B^*\) is integral, \(a\) factors through \(P\).  Thus the
  % maps \(a\:B^*\to P\) and \(b\:X\to B\) are in bijective correspondence.
  % Plainly, this correspondence is compatible with maps \(b'\:B\to B'\) such
  % that \(b'(0)=0\).  In particular, \(1_P\) corresponds to a natural map
  % \(u\:X\to A\) such that \(u(x)=0\), and every map \(b\:X\to B\) factors
  % uniquely through \(u\)."
  \textit{Source: Kleiman, The Picard Scheme, Remark (the Albanese map, rmk:Alb).}
  Let \((C, P)\) be a smooth proper geometrically irreducible curve over \(k\) with a marked \(k\)-rational point \(P \in C(k)\). A smooth proper geometrically irreducible group scheme \(J\) over \(k\) is an \emph{Albanese object} for \((C, P)\) if there exists a morphism \(\iota \colon C \to J\) sending \(P\) to the identity element \(\eta \in J\), and such that every morphism \(f \colon C \to A\) from \(C\) to a smooth proper geometrically irreducible group scheme \(A\) with \(f(P) = \eta_A\) factors uniquely as \(f = \iota \circ g\) for some morphism of group schemes \(g \colon J \to A\).
\end{definition}
\begin{remark}
  \label{rem:IsAlbanese_typeclasses}
  The Lean encoding of Definition~\ref{def:IsAlbanese} takes the four abelian-variety conditions on \(J\) --- group-object structure, properness, smoothness, geometric irreducibility --- as \emph{typeclass parameters} of the declaration \texttt{AlgebraicGeometry.IsAlbanese}, not as conjuncts inside the body. The body is the existential statement on \(\iota\) and its universal-factorisation property; the four hypotheses on \(J\) sit on the binder of the definition. A downstream user constructs an \texttt{IsAlbanese} term against these instances and threads them at every use site (cf.\ the \texttt{letI} chain in \texttt{AbelJacobi.ofCurve} which installs all four instances on \(\Jac(C)\) before applying any Albanese consequence). The same convention is used for the target abelian variety \(A\) in the universal-property clause.
\end{remark}
\subsection{Extracting the universal morphism}
An \texttt{IsAlbanese} term \(h\) packages, behind a single existential, the universal pointed morphism \(\iota \colon C \to J\) together with its two characterising properties: the pointed condition \(P \circ \iota = \eta_J\) and the universal-factorisation property of \(\iota\). The data and the two properties are extracted from the existential body by the following trio of helpers, each implemented in Lean via \texttt{Classical.choose} / \texttt{Classical.choose\_spec} on the existential content of \texttt{IsAlbanese}.
\begin{definition}
\leanok
[The Albanese universal morphism]
  \label{def:IsAlbanese_ofCurve}
  \uses{def:IsAlbanese}

  \lean{AlgebraicGeometry.IsAlbanese.ofCurve}
  Given a pointed curve \((C, P)\) and an Albanese term \(h \colon \mathrm{IsAlbanese}\,C\,P\,J\), the \emph{universal morphism} of \(h\) is the morphism
  \[
    \iota := h.\mathtt{ofCurve} \colon C \to J,
  \]
  obtained by \texttt{Classical.choose} on the existential body of \(h\).
\end{definition}
\begin{lemma}
\leanok
[Pointed condition for the universal morphism]
  \label{lem:IsAlbanese_comp_ofCurve}
  \uses{def:IsAlbanese, def:IsAlbanese_ofCurve}

  \lean{AlgebraicGeometry.IsAlbanese.comp_ofCurve}
  The universal morphism \(h.\mathtt{ofCurve}\) sends the marked point \(P\) to the identity of \(J\):
  \[
    P \circ h.\mathtt{ofCurve} \;=\; \eta_J.
  \]
\end{lemma}
\begin{lemma}
\leanok
[Universal factorisation through the Albanese object]
  \label{lem:IsAlbanese_exists_unique_ofCurve_comp}
  \uses{def:IsAlbanese, def:IsAlbanese_ofCurve}

  \lean{AlgebraicGeometry.IsAlbanese.exists_unique_ofCurve_comp}
  For every morphism \(f \colon C \to A\) from \(C\) to a smooth proper geometrically irreducible group scheme \(A\) over \(k\) with \(P \circ f = \eta_A\), there exists a unique morphism \(g \colon J \to A\) such that \(f = h.\mathtt{ofCurve} \circ g\).
\end{lemma}
This extraction API is precisely what the protected interface \(\mathtt{Jacobian.ofCurve}\), \(\mathtt{Jacobian.comp\_ofCurve}\), \(\mathtt{Jacobian.exists\_unique\_ofCurve\_comp}\) of Chapter~\ref{chap:AbelJacobi} consumes: each protected wrapper takes a marked point \(P \in C(k)\) as input, projects the per-point Albanese data from the bundled witness \texttt{JacobianWitness} (Definition~\ref{def:JacobianWitness}) at that \(P\), and applies the corresponding helper above.
\begin{theorem}\leanok
[Uniqueness of the Albanese object]
  \label{thm:IsAlbanese_unique}
  \uses{def:IsAlbanese, def:IsAlbanese_ofCurve}
  \lean{AlgebraicGeometry.IsAlbanese.unique}
  Let \(J_1, J_2\) be two Albanese objects for the same pointed curve \((C, P)\), with universal pointed morphisms \(\iota_1 := h_1.\mathtt{ofCurve} \colon C \to J_1\) and \(\iota_2 := h_2.\mathtt{ofCurve} \colon C \to J_2\) (Definition~\ref{def:IsAlbanese_ofCurve}). Then there is a unique morphism \(e \colon J_1 \to J_2\) satisfying \(\iota_2 = \iota_1 \circ e\). The proof is the standard universal-property argument: each of the two universal morphisms factors through the other, the two composites equal the unique self-factorisation \(\mathrm{id}\), and the produced morphism \(e\) is therefore an isomorphism; only the existence and uniqueness of \(e\) as a morphism are exposed in the conclusion (see Remark~\ref{rem:IsAlbanese_unique_iso}).
\end{theorem}
\begin{remark}
  \label{rem:IsAlbanese_unique_iso}
  The Lean proof of Theorem~\ref{thm:IsAlbanese_unique} internally produces an inverse morphism \(h \colon J_2 \to J_1\) and verifies \(g \circ h = \mathrm{id}_{J_1}\) and \(h \circ g = \mathrm{id}_{J_2}\) (the invertibility witnesses are computed as intermediate lemmas of the tactic block) before returning the triple \((g,\ \iota_2 = \iota_1 \circ g,\ \text{uniqueness})\). The invertibility witnesses are computed but not retained in the return type
  \[
    \exists! (e : J_1 \to J_2),\ \iota_2 = \iota_1 \circ e.
  \]
  The natural strengthening would be to return
  \[
    \exists! (e : J_1 \cong J_2),\ \iota_2 = \iota_1 \circ e.\mathtt{hom},
  \]
  packaging the produced morphism together with its inverse. Tightening the conclusion in this way is a Lean-side refactor of the signature of \texttt{IsAlbanese.unique} that the project may carry out in a future iteration. The downstream consumers of \texttt{IsAlbanese.unique} (the protected \texttt{AbelJacobi.Jacobian.*} interface) use only the morphism-and-uniqueness content of the conclusion, so the current weaker conclusion is sufficient for the present API.
\end{remark}
\subsection{The Jacobian definition}
\begin{definition}\leanok
[Jacobian]
  \uses{def:genus, def:IsAlbanese, def:JacobianWitness, thm:nonempty_jacobianWitness}
  \label{def:Jacobian}
  
  \lean{AlgebraicGeometry.Jacobian}
  The \emph{Jacobian} of \(C\) is the abelian variety \(\Jac(C)\) over \(k\) obtained as the Albanese object of \(C\): a smooth proper geometrically irreducible group scheme over \(k\) of dimension \(g(C)\) equipped with a universal morphism \(C \to \Jac(C)\) (for each choice of marked point). It is viewed as an object of the category of \(k\)-schemes over \(\Spec k\). Formally, \(\Jac(C)\) is defined as the underlying scheme of a (uniform-over-\(P\)) Albanese witness for \(C\) provided by Theorem~\ref{thm:nonempty_jacobianWitness}.
  When \(g(C) = 0\) the Jacobian collapses to \(\Spec k\) itself, which is the terminal object in the category of \(k\)-schemes. When \(g(C) > 0\) the underlying scheme is the abelian variety produced by Theorem~\ref{thm:nonempty_jacobianWitness}; the genus-one case recovers \(C\) itself (after choosing a \(k\)-rational point), and the higher genus case yields the non-trivial \(g\)-dimensional abelian variety. Both cases are handled uniformly: \(\Jac(C)\) is the underlying scheme of the Albanese witness, with no special-case branching at the level of the definition. The genus-\(0\) specialisation is implicit in the fact that any group scheme over \(k\) of relative dimension \(0\) is isomorphic to \(\Spec k\).
\end{definition}
\section{Group scheme structure and abelian-variety properties}
The four protected instances on \(\Jac(C)\) are each obtained, at the level of the Lean formalisation, by projecting a field of the underlying Albanese witness produced by Theorem~\ref{thm:nonempty_jacobianWitness}. The mathematical content of each property is the corresponding statement about the Albanese variety of a smooth proper geometrically irreducible curve.
\begin{theorem}\leanok
[Group object structure]
  \label{thm:Jacobian_grpObj}
  \uses{def:Jacobian, thm:nonempty_jacobianWitness}
  \lean{AlgebraicGeometry.Jacobian.instGrpObj}
  \(\Jac(C)\) carries a canonical group-object structure in the category of \(k\)-schemes over \(\Spec k\).
\end{theorem}
\begin{proof}

\leanok
  By Definition~\ref{def:Jacobian}, \(\Jac(C) = (\mathrm{jacobianWitness}\, C).J\) as a \(k\)-scheme, where \(\mathrm{jacobianWitness}\, C\) is the Albanese witness extracted from Theorem~\ref{thm:nonempty_jacobianWitness} via \texttt{Classical.choice}. The witness carries the group-object structure as its field \(\mathtt{grpObj} : \mathrm{GrpObj}\, J\), so \(\Jac(C)\) inherits it by projection.
  Mathematically, any Albanese variety carries a unique group-object structure compatible with its universal pointed morphism. Indeed, if \(\iota \colon C \to J\) is the universal morphism sending the marked point \(P\) to the identity, then the diagonal \(C \times C \to J \times J\), post-composed with the universal morphism out of \(C \times C\) (factored through \(J\) by Theorem~\ref{thm:nonempty_jacobianWitness}), produces the multiplication; the unique morphism to a one-point pointed scheme produces the inverse; the marked point produces the identity. Uniqueness of the resulting group-object structure is again the universal-property argument of Theorem~\ref{thm:IsAlbanese_unique}.
\end{proof}
\begin{theorem}\leanok
[Smoothness and dimension]
  \label{thm:Jacobian_smooth_genus}
  \uses{def:Jacobian, def:genus, thm:nonempty_jacobianWitness}
  \lean{AlgebraicGeometry.Jacobian.smoothOfRelativeDimension_genus}
  % SOURCE: [Kleiman, The Picard scheme], Cor.~cor:sm (smoothness/dimension over a field) + Cor.~cor:ch0 (char-\(0\) smoothness) + Prp.~prp:P0 (smooth over reduced \(S\)); arXiv:math/0504020, pp.~34, 38  (read from references/kleiman-picard-src/kleiman-picard.tex)
  % SOURCE QUOTE (cor:sm): "Assume \(S\) is the spectrum of a field \(k\).
  % Assume \(\IPic_{X/k}\) exists and represents \(\Pic_{(X/k)\et}\).  Then
  % \(\)\dim \IPic_{X/k}\le \dim\tu H^1(\mc O_{\!X}).\(\) Equality holds if and
  % only if \(\IPic_{X/k}\) is smooth at \(0;\) if so, then \(\IPic_{X/k}\) is
  % smooth of dimension \(\dim\tu H^1(\mc O_{\!X})\) everywhere."
  % SOURCE QUOTE (cor:ch0): "Assume \(S\) is the spectrum of a field \(k\).
  % Assume \(\IPic_{X/k}\) exists and represents \(\Pic_{(X/k)\et}\). If \(k\) is
  % of characteristic \(0\), then \(\IPic_{X/k}\) is smooth of dimension
  % \(\dim\tu H^1(\mc O_{\!X})\) everywhere."
  % SOURCE QUOTE (prp:P0, smoothness clause): "Assume \(\IPic_{X/S}\) exists,
  % and represents \(\Pic_{(X/S)\fppf}\).  For \(s\in S\), assume all the
  % \(\IPicz_{X_s/k_s}\) are smooth of the same dimension.  Then \(\IPic_{X/S}\)
  % has an open group subscheme \(\IPicz_{X/S}\) of finite type whose fibers
  % are the \(\IPicz_{X_s/k_s}\).  Furthermore, if \(S\) is reduced, then
  % \(\IPicz_{X/S}\) is smooth over \(S\)."
  \textit{Source: Kleiman, The Picard Scheme, Cor.~(cor:sm), Cor.~(cor:ch0), Prp.~(prp:P0).}
  \(\Jac(C) \to \Spec k\) is smooth of relative dimension \(g(C)\).
\end{theorem}
\begin{proof}

\leanok
  The Lean formalisation projects the field \(\mathtt{smoothGenus} : \mathrm{SmoothOfRelativeDimension}\,(\genus\, C)\,J.\mathtt{hom}\) of the witness \(\mathrm{jacobianWitness}\, C\); the relative dimension \(g(C)\) is recorded directly in the witness so that no separate dimension calculation is needed on the Lean side.
  Mathematically, the dimension of an Albanese variety of \(C\) equals the genus \(g = g(C) = \dim_k H^1(C, \mathcal O_C)\) (Phase~A, Chapter~\ref{chap:Cohomology_StructureSheafAb}). The tangent space to \(\Jac(C)\) at the identity is canonically identified with \(H^1(C, \mathcal O_C)\) by deformation theory of the Picard functor: an infinitesimal deformation of \(\mathcal O_C\) is classified by an element of \(H^1(C, \mathcal O_C)\), and this identifies the tangent space at the identity of \(\Pic^0_{C/k}\) with \(H^1(C, \mathcal O_C)\). Smoothness of \(\Jac(C)\) over \(k\) follows from the smoothness of any group scheme over a field whose tangent space at the identity is of constant rank equal to the relative dimension --- a consequence of homogeneity of the group action, since translation by any \(k\)-rational point is an automorphism and so the local structure of \(\Jac(C)\) at every \(k\)-rational point is the same as at the identity.
\end{proof}
\begin{theorem}\leanok
[Properness]
  \label{thm:Jacobian_proper}
  \uses{def:Jacobian, thm:nonempty_jacobianWitness}
  \lean{AlgebraicGeometry.Jacobian.instIsProper}
  % SOURCE: [Kleiman, The Picard scheme], Prp.~prp:P0 (\(\IPicz_{X/S}\) closed in \(\IPic_{X/S}\) and proper over \(S\) when fibres are complete); arXiv:math/0504020, p.~38  (read from references/kleiman-picard-src/kleiman-picard.tex)
  % SOURCE QUOTE: "Assume \(\IPic_{X/S}\) exists, and represents
  % \(\Pic_{(X/S)\fppf}\).  For \(s\in S\), assume all the \(\IPicz_{X_s/k_s}\) are
  % smooth of the same dimension.  Then \(\IPic_{X/S}\) has an open group
  % subscheme \(\IPicz_{X/S}\) of finite type whose fibers are the
  % \(\IPicz_{X_s/k_s}\).  Furthermore, if \(S\) is reduced, then \(\IPicz_{X/S}\)
  % is smooth over \(S\).  Moreover, if all the \(\IPicz_{X_s/k_s}\) are complete
  % and if \(\IPic_{X/S}\) is separated over \(S\), then \(\IPicz_{X/S}\) is closed
  % in \(\IPic_{X/S}\) and proper over \(S\)."
  \textit{Source: Kleiman, The Picard Scheme, Prp.~(prp:P0).}
  \(\Jac(C) \to \Spec k\) is proper.
\end{theorem}
\begin{proof}

\leanok
  The Lean formalisation projects the field \(\mathtt{proper} : \mathrm{IsProper}\, J.\mathtt{hom}\) of the witness.
  Mathematically, properness of the Albanese variety is part of its construction. In the Picard-scheme route, \(\Pic^0_{C/k}\) is a closed subgroup scheme of the relative Picard scheme \(\Pic_{C/k}\), which is itself proper over \(\Spec k\) because \(C\) is proper and geometrically connected (FGA representability for \(\Pic\) of a smooth proper geometrically connected curve produces a proper group scheme; cf.\ FGA Explained, Chapter~9). In the symmetric-power route, properness of \(\Pic^g_{C/k}\) (and hence \(\Pic^0_{C/k}\) after translation) follows from properness of \(\Sym^g(C)\) together with the universal closedness of the surjection \(\sigma\). In either case, properness descends to \(\Jac(C)\).
\end{proof}
\begin{theorem}\leanok
[Geometric irreducibility]
  \label{thm:Jacobian_geomIrred}
  \uses{def:Jacobian, thm:nonempty_jacobianWitness}
  \lean{AlgebraicGeometry.Jacobian.instGeometricallyIrreducible}
  % SOURCE: [Kleiman, The Picard scheme], Lem.~lem:agps (the identity component \(G^0\) is open-closed, finite type, geometrically irreducible, commutes with field extension) + Prp.~prp:pic0 (specialisation to \(\IPicz_{X/k}\)); arXiv:math/0504020, pp.~30, 31  (read from references/kleiman-picard-src/kleiman-picard.tex)
  % SOURCE QUOTE (lem:agps): "Let \(k\) be a field, and \(G\) a group scheme
  % locally of finite type.  Let \(G^0\) denote the connected component of the
  % identity element \(e\).  \tu{(1)} Then \(G\) is separated.  \tu{(2)} Then \(G\)
  % is smooth if it has a geometrically reduced open subscheme.  \tu{(3)}
  % Then \(G^0\) is an open and closed group subscheme of finite type; it is
  % geometrically irreducible; and forming it commutes with extending \(k\)."
  % SOURCE QUOTE (prp:pic0): "Assume \(S\) is the spectrum of a field \(k\).
  % Assume \(\IPic_{X/k}\) exists and represents \(\Pic_{(X/k)\fppf}\).  Then
  % \(\IPic_{X/k}\) is separated, and it is smooth if it has a geometrically
  % reduced open subscheme.  Furthermore, the connected component of the
  % identity \(\IPicz_{X/k}\) is an open and closed group subscheme of finite
  % type; it is geometrically irreducible; and forming it commutes with
  % extending \(k\)."
  \textit{Source: Kleiman, The Picard Scheme, Lem.~(lem:agps), Prp.~(prp:pic0).}
  \(\Jac(C) \to \Spec k\) is geometrically irreducible.
\end{theorem}
\begin{proof}

\leanok
  The Lean formalisation projects the field \(\mathtt{geomIrred} : \mathrm{GeometricallyIrreducible}\, J.\mathtt{hom}\) of the witness.
  Mathematically, the Albanese variety \(\Jac(C)\) is the identity component of \(\Pic_{C/k}\) (the Picard scheme), and the identity component of a group scheme of finite type over a field is geometrically irreducible (in fact, geometrically connected, and since it is smooth at the identity by Theorem~\ref{thm:Jacobian_smooth_genus} it is geometrically irreducible). Equivalently, in the symmetric-power route, \(\Jac(C)\) arises as the Stein factorisation of \(\Sym^g(C) \to \Pic^g_{C/k}\) followed by a translation; \(\Sym^g(C)\) is geometrically irreducible (being the quotient of \(C^g\), itself geometrically irreducible as a product of geometrically irreducible factors, by a finite group action), so its image under the Stein factorisation is geometrically irreducible as well.
\end{proof}
\subsection{Sanity check in low genus}
\begin{itemize}
  \item \(g = 0\): Then \(\Jac(C) = \Spec k\), the trivial group scheme; \(\Jac(\mathbb P^1_k) = \Spec k\) is smooth of relative dimension \(0 = g\).
  \item \(g = 1\): Then \(\Jac(C) \cong C\) once a \(k\)-rational point on \(C\) is fixed; \(C\) is itself an elliptic curve (abelian variety of dimension \(1\)).
  \item \(g \ge 2\): \(\Jac(C)\) is a non-trivial \(g\)-dimensional abelian variety, never isomorphic to \(C\) as a scheme.
\end{itemize}
\section{Existence of an Albanese variety}
The classical construction of the Jacobian goes through the Picard scheme machinery: \(\Jac(C)\) is the connected component of the identity of \(\Pic^0_{C/k}\). An equivalent route, available whenever a \(k\)-rational point \(P \in C(k)\) has been fixed, is to take the Stein factorisation of the Abel--Jacobi morphism \(\Sym^g(C) \to \Pic^g_{C/k}\) and translate by \(-gP\).
\subsection{The Albanese witness bundle}
Before stating the existence theorem, we package its conclusion into a single structure --- the \emph{Albanese witness} of \(C\). This bundling reverses the quantifier order of the classical statement (see Remark~\ref{rem:JacobianWitness_quantifier_order}) and is the data structure consumed by the \texttt{Jacobian} family of declarations in this file and by the \texttt{AbelJacobi} interface of Chapter~\ref{chap:AbelJacobi}.
\begin{definition}
\leanok
  \uses{def:IsAlbanese, def:genus}
[Albanese witness for a curve]
  \label{def:JacobianWitness}
  
  \lean{AlgebraicGeometry.JacobianWitness}
  An \emph{Albanese witness} for the smooth proper geometrically irreducible curve \(C\) over \(k\) is a bundle of the following seven fields:
  \begin{itemize}
    \item \(J\): an object of the category of \(k\)-schemes over \(\Spec k\) --- the candidate Albanese object.
    \item \(\mathtt{grpObj}\): a group-object structure on \(J\) in the category of \(k\)-schemes over \(\Spec k\).
    \item \(\mathtt{proper}\): a proof that the structure morphism \(J \to \Spec k\) is proper.
    \item \(\mathtt{smooth}\): a proof that the structure morphism \(J \to \Spec k\) is smooth.
    \item \(\mathtt{geomIrred}\): a proof that the structure morphism \(J \to \Spec k\) is geometrically irreducible.
    \item \(\mathtt{smoothGenus}\): a proof that the structure morphism \(J \to \Spec k\) is smooth of relative dimension \(g(C)\). This refines \(\mathtt{smooth}\) at the specific relative dimension \(\genus\, C\) (see Remark~\ref{rem:JacobianWitness_smooth_redundancy}).
    \item \(\mathtt{isAlbaneseFor}\): a proof that, for every \(k\)-rational marked point \(P \in C(k)\), the scheme \(J\) together with its group-object, proper, smooth, and geometrically-irreducible structures satisfies the Albanese universal property for the pointed curve \((C, P)\) in the sense of Definition~\ref{def:IsAlbanese}.
  \end{itemize}
  Formally, an Albanese witness is a Lean-level record (a structure) whose name is \texttt{JacobianWitness}.
\end{definition}
\begin{remark}
  \label{rem:JacobianWitness_quantifier_order}
  The classical phrasing of the existence statement of an Albanese variety is
  \[
    \forall\, P \in C(k),\ \exists\, J,\ J \text{ is an Albanese object for } (C, P).
  \]
  The witness bundle reverses these quantifiers: a single \(J\) is chosen once (intrinsic to \(C\)), and the marked-point-dependent data is recorded as the family-valued field
  \[
    \mathtt{isAlbaneseFor} \colon \forall\, P \in C(k),\ \mathrm{IsAlbanese}\,C\,P\,J.
  \]
  Mathematically this reversal is justified because the underlying scheme of any Albanese object is intrinsic to \(C\): a different choice of \(P \in C(k)\) alters the universal morphism \(\iota_P \colon C \to J\) by a translation of \(J\), which is an automorphism of \(J\), and so leaves \(J\) unchanged as a scheme. The reversed bundling is what makes the per-\(P\) projection in Chapter~\ref{chap:AbelJacobi} clean: \(\mathtt{Jacobian.ofCurve}\), \(\mathtt{Jacobian.comp\_ofCurve}\), and \(\mathtt{Jacobian.exists\_unique\_ofCurve\_comp}\) each take a marked point \(P\) as input and project the corresponding component of \(\mathtt{isAlbaneseFor}\,P\) via the extraction trio (Definition~\ref{def:IsAlbanese_ofCurve}, Lemmas~\ref{lem:IsAlbanese_comp_ofCurve}, \ref{lem:IsAlbanese_exists_unique_ofCurve_comp}).
\end{remark}
\begin{remark}
  \label{rem:JacobianWitness_smooth_redundancy}
  The fields \(\mathtt{smooth}\) and \(\mathtt{smoothGenus}\) of an Albanese witness are mathematically redundant: the relative-dimension version implies the unqualified version via Mathlib's \texttt{AlgebraicGeometry.SmoothOfRelativeDimension.smooth}. Keeping both fields is a Lean-level convenience: \(\mathtt{smooth}\) is the form required by the typeclass binder of \texttt{IsAlbanese} (and by the protected \(\mathtt{Jacobian.instGrpObj}\) projection), whereas \(\mathtt{smoothGenus}\) is the form required by the protected \(\mathtt{Jacobian.smoothOfRelativeDimension\_genus}\) projection. The redundancy carries no mathematical cost (the two fields agree) and no soundness cost (deriving one from the other would only re-run the same Mathlib instance synthesis at the projection sites).
  \uses{def:IsAlbanese, def:genus, def:JacobianWitness}
\end{remark}
\begin{theorem}\leanok
[Existence of an Albanese variety uniformly over the marked point]
  \label{thm:nonempty_jacobianWitness}
  \uses{def:JacobianWitness, def:picardJacobianWitness}

  \lean{AlgebraicGeometry.nonempty_jacobianWitness}
  % Cycle-fix (iter-149 plan-agent): drop the spurious 
  % from this statement-block. Mathematically `def:Jacobian` is the entity
  % that depends on this existence theorem (def:Jacobian's  lists
  % thm:nonempty_jacobianWitness), not the other way around. Listing
  % def:Jacobian here as a hypothesis of the existence theorem introduced
  % a 2-cycle in plastex's depgraph (def:Jacobian ↔ thm:nonempty_jacobianWitness)
  % that crashes `leanblueprint web` with RecursionError in depgraph.ancestors.
  % SOURCE: [Kleiman, The Picard scheme], Exercise ex:jac (generalized Jacobians of a family of curves form a smooth quasi-projective family of relative dimension \(p_a\), projective iff \(X/S\) smooth); arXiv:math/0504020, p.~50  (read from references/kleiman-picard-src/kleiman-picard.tex)
  % SOURCE QUOTE: "Assume \(X/S\) is projective and flat, its geometric fibers
  % are integral curves of arithmetic genus \(p_a\), and \(S\) is Noetherian.
  % Show the ``generalized Jacobians'' \(\IPicz_{X_s/k_s}\) form a smooth
  % quasi-projective family of relative dimension \(p_a\).  And show this
  % family is projective if and only if \(X/S\) is smooth."
  \textit{Source: Kleiman, The Picard Scheme, Exercise (ex:jac).}
  Let \(C\) be a smooth proper geometrically irreducible curve over a field \(k\). There exists a smooth proper geometrically irreducible group scheme \(J\) over \(k\) of relative dimension \(g(C)\) such that, for every \(k\)-rational point \(P \in C(k)\), the scheme \(J\) is an Albanese object for the pointed curve \((C, P)\) in the sense of Definition~\ref{def:IsAlbanese}. Equivalently, the type of Albanese witnesses for \(C\) (Definition~\ref{def:JacobianWitness}) is non-empty.
\end{theorem}
\begin{proof}

\leanok
  The witness is constructed uniformly as the Picard identity component \(J = \Pic^0_{C/k}\) for every genus; the genus-\(0\) case requires no separate treatment, since there \(\Pic^0_{C/k} = \Spec k\) automatically. The Albanese universal property of \(J\) consumes the general Milne \S I.1--I.3 rigidity-lemma corollaries of \cref{chap:AbelianVarietyRigidity}. The route below produces the intrinsic scheme \(J\); only the universal morphism \(\iota_P \colon C \to J\) varies with the marked point \(P\). We outline the construction and name the Mathlib infrastructure it requires.
  \paragraph{(0) Reduction to a \(P\)-independent witness.}
  The witness \(J\) depends only on \(C\), not on the marked point \(P\). Indeed, both classical constructions below produce \(J\) as an intrinsic object of \(C\) (the identity component of \(\Pic_{C/k}\), or the Stein-factorisation target of \(\Sym^g(C) \to \Pic^g_{C/k}\) translated by a fixed \(P_0\)); a different choice of \(P\) in \(C(k)\) alters \(\iota_P\) by a translation of \(J\), which is an automorphism, and so does not change \(J\) as a scheme. The Lean encoding reflects this by bundling \(J\) as a single field and the universal-property data as \(\forall P,\, \mathrm{IsAlbanese}\, C\, P\, J\). We now turn to producing \(J\).
  \paragraph{Route A --- Picard scheme.}
  This route follows Hartshorne~III.4 and FGA Explained, Chapter~9.
  \begin{itemize}
    \item[\textbf{(A.1)}] \emph{The relative Picard functor.} For a \(k\)-scheme \(T\), write \(C_T := C \times_k T\) and let \(\pi \colon C_T \to T\) be the projection. The relative Picard functor is
    \[
      \Pic^{\sharp}_{C/k} \colon (\Sch/k)^{\mathrm{op}} \to \Sets, \qquad
      T \longmapsto \Pic(C_T) / \pi^*\Pic(T).
    \]
    A \(T\)-point is a line bundle on \(C_T\) modulo a line bundle pulled back from \(T\). The set-valued formulation suffices for representability; the canonical abelian-group structure on \(\Pic(C_T)\) descends to the quotient by \(\pi^*\Pic(T)\) and gives \(\Pic^{\sharp}_{C/k}\) a natural group-functor refinement.
    \item[\textbf{(A.2)}] \emph{Representability.} For \(C\) smooth proper geometrically connected over \(k\) with a \(k\)-rational point, \(\Pic^{\sharp}_{C/k}\) (after \'etale sheafification, if \(k\) is not algebraically closed) is representable by a \(k\)-group scheme \(\Pic_{C/k}\), locally of finite type. This is FGA's main representability theorem (Grothendieck, FGA~232); the smooth-proper-geometrically-connected hypothesis ensures that the cohomology-and-base-change machinery applies and that \(H^0(C, \mathcal O_C) = k\), which is needed for the quotient by \(\pi^*\Pic(T)\) to be \'etale-locally injective.
    \item[\textbf{(A.3)}] \emph{The identity component.} Let \(\Pic^0_{C/k}\) be the identity component of \(\Pic_{C/k}\). As the identity component of a \(k\)-group scheme locally of finite type, \(\Pic^0_{C/k}\) is geometrically irreducible; for \(C\) smooth proper of genus \(g\), it is moreover a smooth proper \(k\)-group scheme of dimension \(g\). Smoothness comes from the deformation-theoretic identification of its tangent space at the identity with \(H^1(C, \mathcal O_C)\) and homogeneity of the group action; properness from the fact that the natural map \(\Pic^0_{C/k} \to \Pic_{C/k}\) is an open and closed immersion onto a proper component (a consequence of the cohomological criterion: \(\Pic^0_{C/k}\) is the kernel of the degree map \(\Pic_{C/k} \to \underline{\mathbb Z}_k\), which is locally constant). Take \(J := \Pic^0_{C/k}\).
    \item[\textbf{(A.4)}] \emph{The Abel--Jacobi universal morphism.} For each \(P \in C(k)\), the Abel--Jacobi morphism
    \[
      \iota_P \colon C \to \Pic^0_{C/k}, \qquad
      Q \longmapsto [\mathcal O_C(Q - P)]
    \]
    sends \(P\) to the identity of \(\Pic^0_{C/k}\). It is universal among morphisms \(f \colon C \to A\) to smooth proper geometrically irreducible \(k\)-group schemes \(A\) with \(f(P) = \eta_A\): this is the classical Albanese property of \(\Pic^0_{C/k}\), proved by reducing to the case \(k = \overline k\) (Galois descent), where it is a standard consequence of the autoduality of the Jacobian and the fact that line bundles on \(C \times A\) trivialised along \(\{P\} \times A\) correspond bijectively to morphisms \(C \to A^{\vee} = A\) sending \(P\) to \(\eta_A\).
  \end{itemize}
  \emph{Mathlib status for Route A.} Each sub-step is missing.
  \begin{itemize}
    \item Sub-step A.1 requires a working theory of relative line bundles on a product \(C \times_k T\) and the descent of the abelian-group structure on \(\Pic\) along \(\pi^*\). The pieces are present individually but not assembled. (See per-sub-phase budget below: \(\sim 700\)--\(1100\) LOC, \(\sim 6\)--\(10\) iters.)
    \item Sub-step A.2 requires FGA representability for \(\Pic_{C/k}\), which in turn requires Hilbert / Quot scheme infrastructure not in Mathlib (\texttt{HilbertScheme}, \texttt{QuotScheme}, and the cohomology-and-base-change for proper morphisms with finitely presented kernels). The construction engine for this Quot/Hilbert layer is Nitsure, ``Construction of Hilbert and Quot Schemes'' (FGA Explained; arXiv:math/0504590), \S\,4 (generic flatness and the flattening stratification) and \S\,5 (the construction of \(\mathrm{Quot}\), hence \(\mathrm{Hilb}\), as a scheme); see \texttt{references/nitsure-hilbert-quot.md}. (See per-sub-phase budget below: \(\sim 2200\)--\(3000\) LOC, \(\sim 15\)--\(25\) iters --- the dominant Route A block.)
    \item Sub-step A.3 requires connectedness machinery for group schemes locally of finite type, the identity-component construction \(G^0 \subseteq G\), and the degree map \(\Pic_{C/k} \to \underline{\mathbb Z}_k\). (See per-sub-phase budget below: \(\sim 600\)--\(900\) LOC, \(\sim 5\)--\(8\) iters.)
    \item Sub-step A.4 requires the universal property of \(\Pic^0_{C/k}\), which is the classical Albanese functoriality and is not in Mathlib in this form. \textbf{Bypass via Milne autoduality FAILS}: as Milne III \S~6 writes the proof of Proposition~6.1, it invokes Theorem~3.2 (rational-map-extension from a nonsingular variety to an abelian variety) directly to upgrade the symmetric-product rational map \(C^{(g)} \dashrightarrow A\) to a regular map \(J \to A\), and the would-be Picard-functor / autoduality detour \(J \cong J^\vee\) routes through the theorem of the cube, excised from this project iter-163. Sub-step~A.4 therefore inherits the Theorem~3.2 + Lemma~3.3 codim-\(1\) extension + Weil-divisor API + Auslander--Buchsbaum sub-build (the last absent from Mathlib). (See iter-172 audit at \cref{sec:Jacobian_routeA4_albaneseUP}, NOTE block + \cref{thm:albanese_universal_property_northstar} proof.) (See per-sub-phase budget below: \(\sim 2500{+}\) LOC, \(\sim 22\)--\(35\) iters.)
  \end{itemize}
  \paragraph{Route A --- per-sub-phase LOC and iter budget.}
  The Route A build-out breaks into four sub-phases with the following cost
  estimates (LOC = lines of Lean to land axiom-clean; iters = prover iterations
  at the project's current \(\sim 80\)--\(100\) net-LOC-per-iter velocity on
  infrastructure material). The Mathlib-prerequisite cascade for each
  sub-phase is also recorded so subsequent planner decisions can sequence the
  upstream pieces.
  \begin{itemize}
    \item[(A.1)] \emph{Relative Picard functor + relative-line-bundles
      theory.} Net new project material: \(\sim 700\)--\(1100\) LOC. Iters:
      \(\sim 6\)--\(10\). The smallest standalone Mathlib entry point is the
      \texttt{RelativeSpec} functor (per the iter-123 audit
      \texttt{analogies/m3-route-audit.md}); the rest is line-bundle / \(\Pic\)
      machinery on a product \(C \times_k T\), which Mathlib has only partially
      in the form needed.
    \item[(A.2)] \emph{Hilbert / Quot scheme representability + FGA
      representability of \(\Pic_{C/k}\).} Net new project material:
      \(\sim 2200\)--\(3000\) LOC. Iters: \(\sim 15\)--\(25\). The dominant cost.
      The Quot/Hilbert construction engine is Nitsure \S 4--5
      (\texttt{references/nitsure-hilbert-quot.md}); the cohomology-and-base-change
      layer underneath is Stacks tag 02KH
      (\texttt{references/stacks-coherent.md}). The \'etale-sheafification
      step when \(k\) is not algebraically closed adds \(\sim 200\) LOC and
      requires Mathlib's \'etale-sheafification machinery (verified present:
      \texttt{Mathlib/CategoryTheory/Sites/Sheafification.lean} family,
      together with \texttt{Mathlib/AlgebraicGeometry/Sites/Etale.lean}).
    \item[(A.3)] \emph{Identity component \(\Pic^0_{C/k}\) + degree map.}
      Net new project material: \(\sim 600\)--\(900\) LOC. Iters: \(\sim 5\)--\(8\).
      Connectedness machinery for group schemes locally of finite type; the
      identity component construction \(G^0 \subseteq G\); the degree map
      \(\Pic_{C/k} \to \underline{\mathbb Z}_k\) via the locally-constant
      pushforward (cf.\ \texttt{references/stacks-coherent.md}, base change of
      \(R^i f_*\)). Smoothness via the deformation-theoretic identification of
      the tangent space with \(H^1(C, \mathcal O_C)\) uses the cohomology layer
      of (A.2).
    \item[(A.4)] \emph{Albanese universal property of \(\Pic^0_{C/k}\).}
      Net new project material: \(\sim 2500{+}\) LOC. Iters: \(\sim 22\)--\(35\)
      (iter-172 audit re-estimate; see
      \cref{sec:Jacobian_routeA4_albaneseUP}).
      Milne Proposition~6.1 / 6.4 (\texttt{references/abelian-varieties.pdf}).
      \textbf{Bypass via Milne autoduality FAILS}: A.4 inherits the
      Theorem~3.2 + Lemma~3.3 codim-\(1\) extension + Weil-divisor API +
      Auslander--Buchsbaum sub-build (the last absent from Mathlib). Milne's
      proof of Proposition~6.1 as written invokes Theorem~3.2
      (rational-map-extension from a nonsingular variety to an abelian
      variety) directly to upgrade the symmetric-product rational map
      \(C^{(g)} \dashrightarrow A\) to a regular map \(J \to A\); the
      autoduality detour \(J \cong J^\vee\) would route through the theorem
      of the cube, which the project excised iter-163. A.4 reuses the
      proven Rigidity Lemma + Cor 1.2 + Cor 1.5 from the genus-\(0\) stack
      (already in tree, axiom-clean per iter-162) and inherits the
      Theorem~3.2 reserved Lean target
      (\cref{thm:rational_map_to_av_extends}) from
      \cref{chap:AbelianVarietyRigidity}; the codim-\(1\) half (Lemma~3.3,
      pure-codimension-\(1\) indeterminacy on \(\mathbb P^1 \times \mathbb P^1\),
      with no Mathlib Weil-divisor support; see
      \cref{rmk:thm32_codim1_mathlib_gap}) is the dominant LOC driver.
      Additionally, A.4 inherits a symmetric-powers-of-schemes prerequisite
      shared with the historical Route~B. (See iter-172 audit at
      \cref{sec:Jacobian_routeA4_albaneseUP}, NOTE block +
      \cref{thm:albanese_universal_property_northstar} proof.) Char-free.
  \end{itemize}
  Total Route A budget (iter-172 audit re-estimate of A.4 absorbed):
  \(\sim 6000\)--\(7500{+}\) LOC; \(\sim 48\)--\(78\) iters --- summing
  (A.1) \(700\)--\(1100\) LOC / \(6\)--\(10\) iters, (A.2) \(2200\)--\(3000\) LOC /
  \(15\)--\(25\) iters, (A.3) \(600\)--\(900\) LOC / \(5\)--\(8\) iters, and (A.4)
  \(2500{+}\) LOC / \(22\)--\(35\) iters. The previous total \(\sim 4400\)--\(6200\)
  LOC / \(\sim 33\)--\(54\) iters under-counted A.4 (see
  \cref{sec:Jacobian_routeA4_albaneseUP}: the Milne autoduality bypass
  fails and A.4 inherits the Theorem~3.2 + Lemma~3.3 + Auslander--Buchsbaum
  sub-build). The STRATEGY.md \texttt{Iters left}: \texttt{\textasciitilde 40-70}
  envelope is exceeded on the upper end under this re-estimate; the planner
  should treat Route A as a multi-phase build with (A.2) Quot/Hilbert and
  (A.4) Theorem~3.2 codim-\(1\) as the two dominant blocks.
  \paragraph{Mathlib-prerequisite cascade for Route A.}
  The work landings can be partially serialised because each sub-phase has a
  prerequisite tail rooted in a different Mathlib namespace:
  \begin{itemize}
    \item (A.1) \(\to\) \texttt{Mathlib.AlgebraicGeometry.RelativeSpec}
      \texttt{[not yet in Mathlib]},
      \texttt{Mathlib.AlgebraicGeometry.LineBundle.Pullback}
      \texttt{[not yet in Mathlib]},
      \texttt{Mathlib.CategoryTheory.Sites.Pullback} (present).
    \item (A.2) \(\to\) \texttt{Mathlib.AlgebraicGeometry.HilbertScheme}
      (absent),
      \texttt{Mathlib.AlgebraicGeometry.QuotScheme} (absent),
      \texttt{Mathlib.AlgebraicGeometry.Cohomology.BaseChange} (Stacks 02KH;
      partially in Mathlib --- the base-change predicate exists across the
      \texttt{Mathlib/AlgebraicGeometry/Morphisms/} family, but the
      cohomology-and-base-change theorem itself is absent).
    \item (A.3) \(\to\)
      \texttt{Mathlib.AlgebraicGeometry.GroupScheme.IdentityComponent}
      \texttt{[not yet in Mathlib]},
      \texttt{Mathlib.AlgebraicGeometry.LocallyConstantPushforward}
      \texttt{[not yet in Mathlib]}.
      The closest extant anchors are
      \texttt{Mathlib/AlgebraicGeometry/Group/Abelian.lean} and
      \texttt{Mathlib/AlgebraicGeometry/Group/Smooth.lean}.
    \item (A.4) \(\to\) Theorem~3.2 (\cref{thm:rational_map_to_av_extends})
      + Lemma~3.3 codim-\(1\) extension + Weil-divisor API (some shared with
      RR.1) + Auslander--Buchsbaum (absent from Mathlib); also reuses
      \texttt{AlgebraicJacobian.AbelianVarietyRigidity} (in tree,
      axiom-clean) for the Rigidity Lemma + Cor 1.2 + Cor 1.5, plus the
      symmetric-powers-of-schemes prerequisite shared with the historical
      Route~B and Mathlib's Albanese-style universal property machinery.
      \textbf{Bypass via Milne autoduality FAILS}: the Picard-functor /
      \(J \cong J^\vee\) detour around Theorem~3.2 routes through the
      theorem of the cube, excised iter-163. (See iter-172 audit at
      \cref{sec:Jacobian_routeA4_albaneseUP}, NOTE block +
      \cref{thm:albanese_universal_property_northstar} proof.)
  \end{itemize}
  The dominant block is (A.2). It can be started in parallel with the
  genus-\(0\) stack because it is file-disjoint and its prerequisites are
  Mathlib-side rather than project-side.
  \paragraph{Route B --- Symmetric powers and Stein factorisation (historical alternative; not pursued).}
  \textbf{Iter-144 disposition.} Route B is a historical alternative \emph{not pursued} by the project. It is preserved here for scholarly context only. Route A --- Picard scheme via FGA --- is the committed M3 route per STRATEGY.md \S~M3 iter-144 disposition. The decomposition, gating Mathlib pieces, and midpoint LOC figure (\(\sim 9000\) LOC) recorded below are the iter-123 audit values (\texttt{analogies/m3-route-audit.md}); the chapter does not re-cost Route B under current Mathlib snapshots.
  This route follows Milne's \emph{Abelian Varieties} (Chapter~III) and avoids FGA representability at the cost of requiring quotients of schemes by finite group actions.
  \begin{itemize}
    \item[\textbf{(B.1)}] \emph{The \(g\)-fold symmetric power.} Form the product \(C^g\) and let \(S_g\) act by permutation of the factors. The categorical quotient \(C^{(g)} := C^g / S_g\) exists as a \(k\)-scheme, smooth and proper over \(k\) of dimension \(g\). Smoothness follows because the action of \(S_g\) on the smooth \(k\)-scheme \(C^g\) is free outside the diagonal locus and the local-ring computation at the diagonal points produces a smooth ring (this is the symmetric-product computation of Fogarty / Mumford). Properness descends from properness of \(C^g\) via the geometric-quotient property.
    \item[\textbf{(B.2)}] \emph{The Abel--Jacobi morphism on symmetric powers.} Fix a \(k\)-rational point \(P_0 \in C(k)\). The map
    \[
      \sigma \colon C^{(g)} \longrightarrow \Pic^g_{C/k}, \qquad
      \sum_{i=1}^g [Q_i] \longmapsto \mathcal O_C\bigl(\textstyle\sum_{i=1}^g Q_i\bigr)
    \]
    is well-defined (the right-hand side is symmetric in the \(Q_i\) and depends only on the unordered tuple) and is surjective. By Brill--Noether--Riemann--Roch, \(\sigma\) is birational onto its image: a generic divisor of degree \(g\) on \(C\) has \(h^0(\mathcal O_C(D)) = 1\), so the fibre \(\sigma^{-1}([\mathcal O_C(D)])\) is a point at generic divisors; the special fibres are projective spaces \(\mathbb P(h^0 - 1)\) over the locus where \(h^0 > 1\).
    \item[\textbf{(B.3)}] \emph{The Stein factorisation.} Because \(\sigma\) is a proper surjection with geometrically connected fibres (each fibre \(\mathbb P^{h^0-1}_k\) is geometrically irreducible), its Stein factorisation
    \[
      C^{(g)} \xrightarrow{\;\sigma_*\;} \Spec_{\Pic^g_{C/k}}\bigl(\sigma_* \mathcal O_{C^{(g)}}\bigr) \longrightarrow \Pic^g_{C/k}
    \]
    presents the second arrow as a finite morphism whose source is the intermediate scheme. Since the fibres of \(\sigma\) are connected, this intermediate scheme equals \(\Pic^g_{C/k}\) scheme-theoretically (the Stein factorisation is trivial on the second arrow when fibres are connected, by Stein's classical theorem). The middle scheme is therefore \(\Pic^g_{C/k}\) itself, which is smooth proper geometrically irreducible of dimension \(g\). Translating by \(-g P_0\) yields \(\Pic^0_{C/k}\), and we take \(J := \Pic^0_{C/k}\) as in Route A. The universal property is then established as in Sub-step A.4.
  \end{itemize}
  \emph{Mathlib status for Route B.} Each sub-step is blocked by a different missing piece.
  \begin{itemize}
    \item Sub-step B.1 requires symmetric powers of schemes (the construction \(\Sym^g\) for a \(k\)-scheme of finite type) and finite-group-scheme quotients of schemes (the quotient \(C^g / S_g\)). Mathlib has finite-group quotients of types and of various algebraic structures, but not yet a usable scheme-level quotient by a finite group action with the smooth-target property.
    \item Sub-step B.2 requires Brill--Noether--Riemann--Roch, including the relative version for families of curves --- the cohomology-and-base-change machinery for proper morphisms.
    \item Sub-step B.3 requires the Stein factorisation theorem for proper morphisms of locally Noetherian schemes (\(f \colon X \to Y\) proper \(\implies\) \(f_* \mathcal O_X\) is a coherent \(\mathcal O_Y\)-algebra). Mathlib has partial coherence-theoretic infrastructure but does not yet have the cohomology of \(\mathcal O\)-modules of finite type along proper morphisms in the form needed.
  \end{itemize}
  \paragraph{Mathlib infrastructure summary.}
  Per the iter-144 M3 disposition (STRATEGY.md \S~M3), the project commits to Route A and treats Route B as a historical alternative not pursued. The Mathlib build-out gating the uniform construction is therefore:
  \begin{itemize}
    \item[\textbf{(\(\alpha\))}] \emph{Hilbert / Quot scheme infrastructure plus FGA representability (COMMITTED M3 infrastructure --- iter-144 disposition).} This is the committed Mathlib build-out that unlocks Route A, which discharges the witness existence \emph{uniformly in the genus}. It is \textbf{mandatory and unconditional}: the witness object must be a genuine \(\genus C\)-dimensional abelian variety \(\Pic^0_{C/k}\), and no other route in this chapter produces it. The genus-\(0\) case requires no separate treatment --- when \(\genus C = 0\) one has \(H^1(C, \mathcal O_C) = 0\), so \(\Pic^0_{C/k} = \Spec k\) automatically, and the Albanese universal property degenerates to the (then trivial) statement that every pointed morphism \(C \to A\) into an abelian variety is constant, an instance of the general rigidity-lemma corollaries of \cref{chap:AbelianVarietyRigidity}. Route~A is the single highest-cost build-out, since Hilbert/Quot schemes are upstream of much of modern algebraic geometry; but once present, the rest of Route A (identity component, smoothness, properness, universal property --- the last being the Albanese property of \(\Pic^0\), Milne \emph{Abelian Varieties} Proposition~6.1) follows by classical arguments already widely formalised in adjacent settings.
    \item[\textbf{(\(\beta\), historical only --- not pursued).}] \emph{Symmetric powers of schemes plus finite-group-scheme quotients plus Stein factorisation for proper morphisms.} Documented as a historical alternative not pursued by the project (Route B above); see \cref{thm:nonempty_jacobianWitness} Route B for context. This entry is retained for scholarly continuity only; under the iter-144 disposition, (\(\beta\)) is \emph{not} a build-out the project is gating any of its deliverables on, and (\(\beta\)) does not appear on any active critical path.
  \end{itemize}
  In sum: Mathlib currently does not contain the active infrastructure build (\(\alpha\)), and the existence statement is recorded here as the project's single explicit foundational hypothesis. It backs all downstream consequences --- the four protected instances on \(\Jac(C)\) (\cref{thm:Jacobian_grpObj,thm:Jacobian_smooth_genus,thm:Jacobian_proper,thm:Jacobian_geomIrred}), and the Abel--Jacobi map and its universal property of \cref{chap:AbelJacobi} --- and is discharged once (\(\alpha\)) lands. The historical entry (\(\beta\)) is preserved for scholarly context only and is not part of the project's discharge path.
  \paragraph{Body structure.} The Lean body of \cref{thm:nonempty_jacobianWitness} delegates uniformly to the single witness \texttt{AlgebraicGeometry.picardJacobianWitness} (\(J = \Pic^0_{C/k}\)), whose body is a single \texttt{sorry} --- the FGA construction of \(\Pic^0_{C/k}\) as an abelian variety. There is no genus case-split: the genus-\(0\) case is the automatic degenerate instance \(\Pic^0_{C/k} = \Spec k\). The witness existence is therefore conditionally established once that body closes (no Mathlib gap directly on the witness-existence theorem itself, only on the construction of \(\Pic^0_{C/k}\)).
\end{proof}
\subsection{The chosen Albanese witness}
Having established that the type of Albanese witnesses for \(C\) is non-empty (\cref{thm:nonempty_jacobianWitness}), we fix one such witness once and for all; it is this chosen witness whose seven fields supply the underlying scheme of \(\Jac(C)\) together with each of its four abelian-variety structures.
\begin{definition}

\leanok
[Chosen Albanese witness]
  \label{def:jacobianWitness}
  \uses{def:JacobianWitness, thm:nonempty_jacobianWitness}
  \lean{AlgebraicGeometry.jacobianWitness}
  For a smooth proper geometrically irreducible curve \(C\) over \(k\), \(\mathtt{jacobianWitness}\,C \colon \mathtt{JacobianWitness}\,C\) is a choice of Albanese witness for \(C\) (\cref{def:JacobianWitness}), extracted from the non-emptiness \cref{thm:nonempty_jacobianWitness} via \texttt{Classical.choice}. Its fields define the Jacobian \(\Jac(C)\) (\cref{def:Jacobian}) and discharge each of the four protected abelian-variety instances on it (group object, smoothness of relative dimension \(g(C)\), properness, geometric irreducibility).
\end{definition}
\begin{proof}
  Proved directly in Lean.
\end{proof}
\subsection{Route A.4: the Albanese universal property of \(\Pic^0_{C/k}\) (Milne III \S 6 audit)}
\label{sec:Jacobian_routeA4_albaneseUP}
% NOTE (iter-172 audit, blueprint-writer slug `a4-bypass-audit`):
% AUDIT OUTCOME (b) --- the bypass FAILS. Milne III \S 6's proof of the Albanese
% universal property (Prop 6.1, the existence-and-uniqueness statement consumed
% by A.4) is NOT derivable from Picard-functor functoriality + Cor 1.5/Cor 1.2 +
% Mathlib cohomology-and-base-change alone. The text of the proof of (6.1)
% invokes (I 3.2) --- the rational-map-extension theorem from a nonsingular
% variety to an abelian variety --- to upgrade the symmetric-product rational
% map \(C^{(g)} \dashrightarrow A\) to a regular map \(J \to A\). Theorem 3.2 of
% Milne \S I.3 is, in turn, the composite of Theorem 3.1 (valuative extension
% outside codimension \(\geq 2\)) with Lemma 3.3 (the indeterminacy locus into a
% group variety is pure codimension 1) --- and Lemma 3.3 is precisely the
% Weil-divisor / Auslander--Buchsbaum-flavoured sub-build that has NO Mathlib
% support (already recorded in \cref{rmk:thm32_codim1_mathlib_gap} of
% \cref{chap:AbelianVarietyRigidity}, with the in-tree target
% \cref{thm:rational_map_to_av_extends} /
% \texttt{AlgebraicGeometry.rationalMap\_to\_av\_extends} reserved but unproven).
% A "pure Picard-functor" detour would need autoduality \(J \cong J^\vee\) (Milne
% III Thm 6.6), which itself rests on the theorem of the cube --- and the cube
% was EXCISED from this project iter-163. Net conclusion: A.4 incorporates a
% Thm 3.2 sub-build inherited from \cref{chap:AbelianVarietyRigidity}, plus a
% symmetric-powers-of-schemes prerequisite (a Mathlib gap also blocking the
% historical Route B), plus the image-sum-fills-\(J\) dimension fact. The
% STRATEGY.md row "A.4 reuses in-tree Rigidity Lemma + Cor 1.2/1.5; no new
% Mathlib namespace" is materially under-counted and should be re-estimated by
% the planner. One-sentence summary of the proof chain (Milne III \S 6 as
% written): symmetrise \(C^g \to A\), descend to \(C^{(g)} \dashrightarrow A\),
% extend to a regular map \(J \to A\) via Thm 3.2, pointed-to-homomorphism by Cor
% 1.2; the analogous Prop 6.4 (for \(\varphi : C \times C \to A\) vanishing on the
% diagonal) reduces to Prop 6.1 via Cor 1.5 (additive decomposition).
\begin{theorem}
[Albanese universal property of \(\Pic^0_{C/k}\) (Milne III Prop 6.1)]
  \label{thm:albanese_universal_property_northstar}
  
  \lean{AlgebraicGeometry.Scheme.Pic.albaneseUP}
  % SOURCE: [Milne, Abelian Varieties], Proposition 6.1, \S III.6, doc p.~104 (PDF p.~110) (read from references/abelian-varieties.pdf, PDF page 110)
  % SOURCE QUOTE: "PROPOSITION 6.1. Let \(P\) be a \(k\)-rational point on \(C\).
  % The map \(f^P\colon C\to J\) has the following universal property: for any
  % map \(\varphi\colon C\to A\) from \(C\) into an abelian variety sending \(P\)
  % to \(0\), there is a unique homomorphism \(\psi\colon J\to A\) such that
  % \(\varphi=\psi\circ f^P\)."
  % SOURCE QUOTE (Prop 6.4, the symmetric variant invoked through Cor 1.5;
  % read from references/abelian-varieties.pdf, PDF page 111): "PROPOSITION
  % 6.4. Let \(A\) be an abelian variety over \(k\). For any map
  % \(\varphi\colon C\times C\to A\) such that \(\varphi(\Delta)=0\), there is a
  % unique homomorphism \(\psi\colon J\to A\) such that \(\psi\circ F=\varphi\)."
  % SOURCE QUOTE (Remark 6.5, the Albanese characterization; read from
  % references/abelian-varieties.pdf, PDF page 111): "REMARK 6.5. The
  % proposition says that \((A,F)\) is the Albanese variety of \(C\) in the
  % sense of Lang 1959, II 3, p45. Clearly the pairs \((J,f^P)\) and \((J,F)\)
  % are characterized by the universal properties in (6.1) and (6.4)."
  \textit{Source: Milne, Abelian Varieties, III \S 6, Proposition 6.1 (p.~104), with Proposition 6.4 and Remark 6.5 (p.~105).}
  Let \(C\) be a smooth proper geometrically irreducible curve of genus \(g \geq 1\) over \(k\), fix a \(k\)-rational point \(P \in C(k)\), write \(J := \Pic^0_{C/k}\) for the identity component of its Picard scheme, and let \(f^P \colon C \to J\), \(Q \mapsto [\mathcal O_C(Q - P)]\), be the Abel--Jacobi morphism. For any morphism \(\varphi \colon C \to A\) from \(C\) to an abelian variety \(A\) over \(k\) with \(\varphi(P) = \eta_A\), there exists a unique homomorphism \(\psi \colon J \to A\) of group schemes such that \(\varphi = \psi \circ f^P\). Equivalently, \(J\) together with \(f^P\) exhibits \(J\) as the Albanese object of the pointed curve \((C, P)\) in the sense of \cref{def:IsAlbanese}.
\end{theorem}
% SOURCE QUOTE PROOF: "PROOF. Consider the map
% \((P_1,\dots,P_g)\mapsto\sum_i\psi(P_i)\colon C^g\to A\). Clearly this is
% symmetric, and so it factors through \(C^{(g)}\). It therefore defines a
% rational map \(\psi\colon J\to A\), which (I 3.2) shows to be a regular map.
% It is clear from the construction that \(\psi\circ f^P=\varphi\) (note that
% \(f^P\) is the composite of \(Q\mapsto Q+(g-1)P\colon C\to C^{(g)}\) with
% \(f^{(g)}\colon C^{(g)}\to J\)). In particular, \(\psi\) maps \(0\) to \(0\), and
% (I 1.2) shows that it is therefore a homomorphism. If \(\psi'\) is a second
% homomorphism such that \(\psi'\circ f^P=\varphi\), then \(\psi\) and \(\psi'\)
% agree on \(f^P(C)+\dots+f^P(C)\) (\(g\) copies), which is the whole of \(J\)."
% [Milne, Prop 6.1, \S III.6, doc p.~104.]
\begin{proof}
  \textit{Source: Milne, Abelian Varieties, III \S 6, Proposition 6.1.}
  Symmetrise the \(g\)-fold sum: the morphism \((Q_1, \dots, Q_g) \mapsto \sum_i \varphi(Q_i) \colon C^g \to A\) is symmetric in the factors and so factors through the \(g\)-fold symmetric product \(C^{(g)} := C^g / S_g\). Composing with the inverse of the Abel--Jacobi map \(f^{(g)} \colon C^{(g)} \to J\), which is birational by Brill--Noether--Riemann--Roch (its generic fibre is a single point), one obtains a \emph{rational} map \(\psi \colon J \dashrightarrow A\). By Milne Theorem~3.2 (\cref{thm:rational_map_to_av_extends}: a rational map from a nonsingular variety to an abelian variety is everywhere defined --- the composite of Theorem~3.1's everywhere-extension-outside-codimension-\(2\) via the valuative criterion with Lemma~3.3's pure-codimension-\(1\) indeterminacy for maps into a group variety) the rational map \(\psi\) extends uniquely to a regular morphism \(J \to A\). The identity \(\psi \circ f^P = \varphi\) holds by construction once \(f^P\) is identified as the composite \(Q \mapsto Q + (g - 1)P \colon C \to C^{(g)}\) followed by \(f^{(g)} \colon C^{(g)} \to J\). The pointed condition \(\psi(0_J) = \eta_A\) then promotes \(\psi\) to a group-scheme homomorphism by Milne Corollary~1.2 (\cref{lem:av_regular_map_is_hom}: a regular map between abelian varieties sending the identity to the identity is automatically a homomorphism). Uniqueness: any two such homomorphisms agree on \(f^P(C)\) and therefore on the sumset \(f^P(C) + \dots + f^P(C)\) (\(g\) summands), which fills \(J\) because \(f^{(g)} \colon C^{(g)} \to J\) is surjective (this is the dimension count \(\dim C^{(g)} = g = \dim J\) together with properness of \(C^{(g)}\)). The companion statement Proposition~6.4 for \(\varphi \colon C \times C \to A\) vanishing on the diagonal \(\Delta\) reduces to Proposition~6.1 via the additive decomposition Corollary~1.5 (\cref{lem:hom_additivity_over_product}): the additivity decomposition \(\varphi = \varphi_1 \circ p_1 + \varphi_2 \circ p_2\) with \(\varphi_i \colon C \to A\), \(\varphi_i(P) = 0\), combined with \(\varphi|_\Delta = 0\), forces \(\varphi_1 = -\varphi_2\), after which Proposition~6.1 supplies the unique \(\psi\).
  \medskip
  \emph{Dependency audit (iter-172, Outcome (b)).} The chain of dependencies of this proof, classified by their formalisation status in the project's tree, is:
  \begin{itemize}
    \item \emph{In-tree, axiom-clean (no new work):} Corollary~1.5 (\cref{lem:hom_additivity_over_product}, \texttt{AlgebraicGeometry.hom\_additive\_decomp\_of\_rigidity}); Corollary~1.2 (\cref{lem:av_regular_map_is_hom}, \texttt{AlgebraicGeometry.av\_regularMap\_isHom\_of\_zero}); the Rigidity Lemma (\cref{thm:rigidity_lemma}) underlying both.
    \item \emph{Reserved Lean target, body not yet closed --- Mathlib gap on the codimension-\(1\) step:} Milne Theorem~3.2 (\cref{thm:rational_map_to_av_extends}, \texttt{AlgebraicGeometry.rationalMap\_to\_av\_extends}). The everywhere-extension-outside-codimension-\(2\) half (Theorem~3.1) is within Mathlib's reach via \texttt{ValuativeCriterion}; the pure-codimension-\(1\) half (Lemma~3.3, Weil-divisor / Auslander--Buchsbaum-flavoured argument that the indeterminacy locus of a map into a group variety is pure codimension~\(1\)) has \emph{no} Mathlib Weil-divisor API and is the route's riskiest sub-build (see \cref{rmk:thm32_codim1_mathlib_gap}). An open question for the prover: whether a \emph{pointwise} valuative extension --- run at each codimension-\(1\) point's discrete valuation ring separately --- side-steps building Weil-divisor theory entirely.
    \item \emph{Mathlib gap, shared with Route B:} symmetric powers of schemes \(C^{(g)} := C^g / S_g\) (Mathlib has no usable scheme-level quotient by a finite group action with the smooth-target property; the historical Route B was blocked on the same item). The Abel--Jacobi surjection \(f^{(g)} \colon C^{(g)} \to J\) and the birationality \(f^{(g)}\) uses Brill--Noether--Riemann--Roch, which routes through the cohomology layer common to Route~A sub-steps (A.1)--(A.3).
    \item \emph{Route A.3 output, consumed here:} the identity component \(J = \Pic^0_{C/k}\) together with its smoothness of dimension \(g\) and the surjectivity / dimension count for \(f^{(g)}\).
    \item \emph{NOT a dependency of A.4 (despite a textbook detour suggesting it might be):} autoduality \(J \cong J^\vee\) and the theorem of the cube. Milne's proof of Prop 6.1 as written does NOT invoke autoduality (which is established as Theorem 6.6 of the same section, \emph{after} 6.1, via the theorem of the cube). A pure Picard-functor / Poincar\'e-sheaf detour around Theorem~3.2 would route through autoduality and the cube, both of which the project has deliberately excised from the genus-\(0\) path (iter-163; \cref{rmk:cube_not_needed}). A.4 therefore commits to the Theorem~3.2 sub-build rather than the cube.
  \end{itemize}
  Per \cref{rmk:thm32_codim1_mathlib_gap}, the pure-codimension-\(1\) step (Lemma~3.3) requires either a Mathlib Weil-divisor / Auslander--Buchsbaum sub-build (a stand-alone Mathlib-prerequisite cascade in its own right, not currently scheduled) or the pointwise-valuative detour mentioned above. The STRATEGY.md A.4 row claim ``no new Mathlib namespace'' is incorrect with respect to this Lemma~3.3 sub-build; the planner should re-estimate A.4 with the codim-\(1\) extension treated as an explicit upstream piece.
\end{proof}
\subsection{The Picard witness \(J = \Pic^0_{C/k}\)}
\label{sec:picardJacobianWitness}
The witness existence \cref{thm:nonempty_jacobianWitness} is discharged \emph{uniformly in the genus} by the single construction \(J = \Pic^0_{C/k}\) via Route A (Picard scheme via FGA --- Hilbert/Quot infrastructure plus FGA representability). Route A is \textbf{mandatory and unconditional}: the witness object must be a genuine \(\genus C\)-dimensional abelian variety \(\Pic^0_{C/k}\). The genus-\(0\) case requires no separate construction --- when \(\genus C = 0\) one has \(H^1(C, \mathcal O_C) = 0\), so \(\Pic^0_{C/k} = \Spec k\) automatically, and the Albanese universal property degenerates to the (then trivial) statement that every pointed morphism \(C \to A\) into an abelian variety is constant, an instance of the general rigidity-lemma corollaries of \cref{chap:AbelianVarietyRigidity}. It is the committed M3 route per the iter-144 disposition in STRATEGY.md \S~M3. Route B (symmetric powers + Stein factorisation) is a historical alternative not pursued by the project.
\begin{theorem}\leanok
[The Picard Albanese witness \(J = \Pic^0_{C/k}\)]
  \label{def:picardJacobianWitness}
  \uses{def:JacobianWitness, def:genus, thm:fga_pic_representability, thm:albanese_universal_property_northstar}

  \lean{AlgebraicGeometry.picardJacobianWitness}
  % SOURCE: [Kleiman, The Picard scheme], Remark rmk:Jac (the Jacobian: \(J:=\IPicz_{X/S}\) of a smooth projective family of genus-\(g>0\) curves exists and is a projective Abelian \(S\)-scheme); arXiv:math/0504020, p.~50  (read from references/kleiman-picard-src/kleiman-picard.tex)
  % SOURCE QUOTE: "Assume \(X/S\) is projective and smooth, its geometric
  % fibers are connected curves of genus \(g>0\), and \(S\) is Noetherian.  Set
  % \(J:=\IPicz_{X/S}\); it exists and is a projective Abelian \(S\)-scheme by
  % Exercise~REF.  Set \(\wh J:=\IPicz_{J/S}\); it exists, is a
  % projective Abelian scheme, and is ``dual'' to \(J\) by
  % Remark~REF."
  \textit{Source: Kleiman, The Picard Scheme, Remark (the Jacobian, rmk:Jac).}
  Let \(C \in \mathrm{Over}\,(\Spec k)\) be a smooth proper geometrically irreducible curve over \(k\) (of arbitrary genus). Then there exists a \(\mathtt{JacobianWitness}\,C\) whose underlying scheme is the Picard identity component \(\mathrm{Pic}^0_{C/k}\), an abelian variety of dimension \(\genus C\) (when \(\genus C = 0\) this is \(\Spec k\) automatically). The full \(\mathtt{JacobianWitness}\) data (group-object structure, smoothness of relative dimension \(\genus C\), properness, geometric irreducibility, and Albanese universal property uniformly over \(k\)-rational marked points) is supplied by the Route A construction (per \cref{thm:nonempty_jacobianWitness} Route A; Route B is documented as a historical alternative not pursued).
\end{theorem}
\begin{proof}

\leanok
  This is the M3 arm of the witness existence. M3 is \textbf{committed to Route A --- Picard scheme via FGA} (per STRATEGY.md \S~M3 iter-144 disposition; per-sub-phase budget recorded in the \emph{Route A --- per-sub-phase LOC and iter budget} paragraph above, \(\sim 6000\)--\(7500{+}\) LOC total, \(\sim 48\)--\(78\) iters after the iter-172 audit re-estimate of A.4 --- the bypass via Milne autoduality FAILS, so A.4 inherits the Theorem~3.2 + Lemma~3.3 + Auslander--Buchsbaum sub-build at \(\sim 2500{+}\) LOC / \(\sim 22\)--\(35\) iters; the earlier iter-170 refresh estimate \(\sim 4400\)--\(6200\) LOC / \(\sim 33\)--\(54\) iters under-counted A.4, and the iter-123 audit \texttt{analogies/m3-route-audit.md} midpoint was \(\sim 6500\) LOC; the audit-corrected budget breaks out the two dominant blocks as (A.2) at \(\sim 2200\)--\(3000\) LOC and (A.4) at \(\sim 2500{+}\) LOC). The autonomous loop writes this Mathlib material directly in-tree; upstream PR extraction is an OPTIONAL downstream lift, not a project deliverable. The committed route produces the seven fields of \cref{def:JacobianWitness} in positive genus via the following Mathlib-prerequisite path:
  \begin{itemize}
    \item \emph{Route A: Hilbert/Quot infrastructure plus FGA representability of \(\Pic_{C/k}\).} See \cref{thm:nonempty_jacobianWitness} Route A for the four-step decomposition (A.1 relative Picard functor; A.2 representability; A.3 identity component; A.4 Abel--Jacobi universal morphism). The smallest entry point into Route A is the \texttt{RelativeSpec} functor (A.1 prerequisite; \(\sim 700\)--\(1100\) LOC) per the iter-123 audit.
  \end{itemize}
  Route B (symmetric powers \(C^{(g)} := C^g/S_g\) + Stein factorisation, \(\sim 9000\) LOC midpoint per the iter-123 audit) is documented in \cref{thm:nonempty_jacobianWitness} Route B as a historical alternative not pursued by the project.
  M3 sits behind M2 critical-path during the M2 wait window; whether to start Route A scaffolding (smallest entry: \texttt{RelativeSpec} functor at \(\sim 700\)--\(1100\) LOC) in parallel with M2 or strictly after M2 closes is a planner-level scheduling call. Existing iter-150 mathlib-analogist scheduling trigger: dispatched if cumulative M2.body-pile LOC exceeds \(925\) LOC without piece (i.b) closing, OR if M2.a body closure has not landed by iter-160.
  The witness \texttt{AlgebraicGeometry.picardJacobianWitness} is the single uniform body to which \cref{thm:nonempty_jacobianWitness} delegates (there is no genus case-split). Its Lean body is currently \texttt{sorry}; the closure is contingent on Route A of \cref{thm:nonempty_jacobianWitness} being prosecuted in full.
\end{proof}
\section{Implementation route via the Albanese functor}
Because the full Picard-scheme representability theorem is not available in current Mathlib, the formal construction proceeds in two layers:
\begin{enumerate}
  \item \textbf{Layer I --- direct definition.} The Jacobian of \(C\) is defined uniformly as the underlying scheme of the Albanese witness \(J = \Pic^0_{C/k}\) provided by \cref{thm:nonempty_jacobianWitness}, carrying the four properties of \cref{thm:Jacobian_grpObj,thm:Jacobian_smooth_genus,thm:Jacobian_proper,thm:Jacobian_geomIrred} (each obtained by projection from the witness). The genus-\(0\) case is not special: when \(\genus C = 0\) one has \(H^1(C, \mathcal O_C) = 0\), so \(\Pic^0_{C/k} = \Spec k\) automatically, and the Albanese universal property degenerates to the (then trivial) statement that every pointed morphism \(C \to A\) into an abelian variety is constant --- an instance of the general rigidity-lemma corollaries of the upstream \cref{chap:AbelianVarietyRigidity}.
  \item \textbf{Layer II --- Abel--Jacobi compatibility.} \cref{chap:AbelJacobi} introduces the Abel--Jacobi morphism \(C \to \Jac(C)\) and shows that it is universal among morphisms from \(C\) to abelian varieties.
\end{enumerate}
